\chapter{INTRODUCTION}
Følgende informasjon kan tas med her:
- Organisering av rapporten og kort beskrivelse av sammenhengen mellom kapitlene.
- Fullstendig definering av oppgaven med utdypning av hva oppgaven går ut på.
Eksempelvis presiseres mål, problemområde, omfang og avgrensning.
- Henvisning til eventuelle tidligere arbeider om samme emne.
- Målgruppe for rapporten, jfr. mål- og mottakeranalysen.
- Studentenes faglige bakgrunn.
- Prosjektgruppens valgte arbeidsformer, f.eks. samtaler og intervjuer med kilder, reiser,
bedriftsbesøk, litteraturstudier, laboratorieforsøk, modellbygging og testing.
- Terminologibruk. Er eksempelvis faguttrykk brukt i en spesiell betydning?
Innledningen skal ikke inneholde noe om det som er gjort i oppgaven og oppnådde resultater.